% ======================================================================
\chapter[Оператор момента импульса]%
        {Оператор момента импульса. Собственные значения \texorpdfstring{$L_z$ и $L^2$}{Lz и L²}}
\label{t:09}
% ======================================================================

\begin{qbox}
\textbf{Билет 9.} Оператор момента импульса. Собственные функции и собственные значения
оператора $L_z$ ($z$-компоненты момента импульса) в сферических координатах. Собственные значения
оператора квадрата момента импульса $L^2$ (без вывода).
\end{qbox}

\section{Содержание билета}

\textbf{Подтемы:}
\begin{enumerate}
 \item оператор момента импульса $\op{\vect L}=\vect r\times\op{\vect p}$, его компоненты в декартовых координатах;
 \item коммутационные соотношения и их следствие: одновременно измеримы лишь $L^2$ и одна проекция;
 \item $\op L_z=-\ii\hbar\,\partial/\partial\varphi$ в сферических координатах; собственные функции $\ee^{\ii m\varphi}$ и значения $m\hbar$;
 \item собственные значения $L^2=\hbar^2l(l+1)$ \emph{(без вывода)}; сферические функции; векторная модель.
\end{enumerate}

\textbf{Необходимые понятия:} операторы, коммутаторы, собственные функции (\tref{t:05});
сферические координаты $(r,\theta,\varphi)$; классический момент импульса.

\section{Определения и физические принципы}

\begin{definition}[оператор момента импульса]
По правилу соответствия $\vect L=\vect r\times\vect p\to\op{\vect L}=\vect r\times\op{\vect p}=-\ii\hbar\,\vect r\times\nabla$:
\begin{equation}\label{eq:9:L}
  \op L_x=-\ii\hbar\Bigl(y\pdiff{}{z}-z\pdiff{}{y}\Bigr),\quad
  \op L_y=-\ii\hbar\Bigl(z\pdiff{}{x}-x\pdiff{}{z}\Bigr),\quad
  \op L_z=-\ii\hbar\Bigl(x\pdiff{}{y}-y\pdiff{}{x}\Bigr),
\end{equation}
$\op L^2=\op L_x^2+\op L_y^2+\op L_z^2$. Порядок сомножителей несуществен: $y$ и $\op p_z$ коммутируют.
\end{definition}

\begin{definition}[квантовые числа момента]
$l=0,1,2,\dots$ — \emph{орбитальное} (азимутальное) квантовое число, определяющее $L^2=\hbar^2l(l+1)$;
$m=-l,\dots,l$ — \emph{магнитное} квантовое число, $L_z=m\hbar$. Состояния с $l=0,1,2,3,4$
обозначают буквами $s,p,d,f,g$.
\end{definition}

\emph{Физический смысл.} Момент импульса сохраняется в центрально-симметричном поле
($\op L^2$ и $\op L_z$ коммутируют с гамильтонианом); поэтому квантовые числа $l,m$ классифицируют
состояния атомов, ядер, молекул. Квантование момента — следствие однозначности волновой функции
при повороте на $2\pi$, а не постулат, как у Бора.

\emph{Область применимости.} Речь идёт об \emph{орбитальном} моменте частицы, связанном с её
движением в пространстве. Собственный момент (спин, \tref{t:12}) описывается теми же коммутационными
соотношениями, но может принимать и полуцелые значения; орбитальный — только целые.

\section{Основная теория}

\subsection{Коммутационные соотношения}

Из $[\op x,\op p_x]=\ii\hbar$ следует (вывод~\ref{der:9:comm})
\begin{equation}\label{eq:9:comm}
  [\op L_x,\op L_y]=\ii\hbar\op L_z,\qquad [\op L_y,\op L_z]=\ii\hbar\op L_x,\qquad [\op L_z,\op L_x]=\ii\hbar\op L_y,
\end{equation}
\begin{equation}\label{eq:9:L2comm}
  [\op L^2,\op L_x]=[\op L^2,\op L_y]=[\op L^2,\op L_z]=0 .
\end{equation}
\textbf{Следствия.} (1) Разные проекции момента не измеримы одновременно: если $L_z$ определена,
то $L_x$ и $L_y$ не имеют определённых значений (исключение — состояние $l=0$, где все проекции нули).
(2) Квадрат момента и одна (любая) проекция измеримы одновременно. Принято выбирать ось $z$;
состояние задаётся парой $(l,m)$. (3) Вектор $\vect L$ не может быть определён полностью — нельзя
сказать, «куда он направлен».

\subsection{Оператор $\op L_z$ в сферических координатах}

В сферических координатах $x=r\sin\theta\cos\varphi$, $y=r\sin\theta\sin\varphi$, $z=r\cos\theta$
(вывод~\ref{der:9:Lz}):
\begin{equation}\label{eq:9:Lzsph}
  \op L_z=-\ii\hbar\pdiff{}{\varphi}.
\end{equation}
$\op L_z$ действует только на азимутальный угол — это генератор поворотов вокруг оси $z$, подобно тому
как $\op p_x=-\ii\hbar\,\partial/\partial x$ — генератор сдвигов вдоль $x$.

\subsection{Собственные функции и значения $\op L_z$}

Уравнение $\op L_z\Phi=L_z\Phi$ решается элементарно (вывод~\ref{der:9:eig}):
\begin{equation}\label{eq:9:Phim}
  \Phi_m(\varphi)=\frac{1}{\sqrt{2\pi}}\ee^{\ii m\varphi},\qquad L_z=m\hbar,\qquad m=0,\pm1,\pm2,\dots
\end{equation}
Дискретность — следствие требования однозначности: точкам $\varphi$ и $\varphi+2\pi$ отвечает одна и та же
точка пространства. Функции ортонормированы на $[0,2\pi)$:
$\int_0^{2\pi}\Phi_{m'}^*\Phi_m\dd\varphi=\delta_{mm'}$. В состоянии $\Phi_m$ плотность вероятности
$|\Phi_m|^2=1/2\pi$ не зависит от $\varphi$ — азимутальный угол полностью неопределён (аналог
соотношения неопределённостей для пары $\varphi$, $L_z$).

\subsection{Квадрат момента (без вывода)}

В сферических координатах $\op L^2$ действует только на углы:
\begin{equation}\label{eq:9:L2}
  \op L^2=-\hbar^2\Bigl[\frac{1}{\sin\theta}\pdiff{}{\theta}\Bigl(\sin\theta\pdiff{}{\theta}\Bigr)
  +\frac{1}{\sin^2\theta}\pdiff{^2}{\varphi^2}\Bigr]\equiv-\hbar^2\Delta_{\theta,\varphi},
\end{equation}
где $\Delta_{\theta,\varphi}$ — угловая часть лапласиана: $\nabla^2=\frac{1}{r^2}\pdiff{}{r}\bigl(r^2\pdiff{}{r}\bigr)+\frac{1}{r^2}\Delta_{\theta,\varphi}$.

\begin{principle}{Собственные значения $\op L^2$ (без вывода)}
Однозначные и конечные на сфере решения уравнения $\op L^2Y=L^2Y$ существуют только при
\begin{equation}\label{eq:9:L2eig}
  L^2=\hbar^2l(l+1),\qquad l=0,1,2,\dots,
\end{equation}
причём при данном $l$ проекция $L_z=m\hbar$ может принимать $2l+1$ значений $m=-l,-l+1,\dots,l$.
Общие собственные функции $\op L^2$ и $\op L_z$ — \emph{сферические функции}
$Y_{lm}(\theta,\varphi)=\Theta_{lm}(\theta)\,\ee^{\ii m\varphi}$; уровень $L^2$ вырожден с кратностью $2l+1$.
\end{principle}

Примеры нормированных сферических функций ($\int|Y_{lm}|^2\sin\theta\,\dd\theta\,\dd\varphi=1$):
\[
  Y_{00}=\frac{1}{\sqrt{4\pi}},\qquad
  Y_{10}=\sqrt{\frac{3}{4\pi}}\cos\theta,\qquad
  Y_{1,\pm1}=\mp\sqrt{\frac{3}{8\pi}}\sin\theta\,\ee^{\pm\ii\varphi}.
\]
Проверка для $Y_{10}$ дана в выводе~\ref{der:9:Y10}.

\paragraph{Почему $|m|\le l$.} В состоянии с определёнными $l,m$
$\avg{L_x^2}+\avg{L_y^2}=\avg{L^2}-m^2\hbar^2=\hbar^2[l(l+1)-m^2]\ge0$, так как средние квадратов
эрмитовых операторов неотрицательны. Значит, $m^2\le l(l+1)$, т.\,е. $|m|\le l$. При $|m|=l$ всё равно
$\avg{L_x^2}+\avg{L_y^2}=\hbar^2l>0$: момент не может быть направлен точно вдоль оси $z$
(иначе все три проекции были бы определены).

\paragraph{Правдоподобность $l(l+1)$.} Для изотропного (неполяризованного) набора состояний
с данным $l$ все три оси равноправны, поэтому $\avg{L^2}=3\avg{L_z^2}$. Усредняя $m^2$ по $2l+1$
равновероятным значениям:
\[
  \avg{L_z^2}=\frac{\hbar^2}{2l+1}\sum_{m=-l}^{l}m^2=\frac{\hbar^2}{2l+1}\cdot\frac{l(l+1)(2l+1)}{3}=\frac{\hbar^2l(l+1)}{3},
\]
откуда $\avg{L^2}=\hbar^2l(l+1)$. Это не вывод~\eqref{eq:9:L2eig}, но наглядное объяснение,
почему $L^2$ больше $(l\hbar)^2$.

\subsection{Векторная модель}

Наглядно вектор $\vect L$ длины $\hbar\sqrt{l(l+1)}$ прецессирует вокруг оси $z$ по конусу так, что
проекция $L_z=m\hbar$ фиксирована, а $L_x$, $L_y$ не определены (рис.~\ref{fig:9:vector}). Угол
конуса $\cos\theta_m=m/\sqrt{l(l+1)}$ никогда не равен 0 или $\pi$. При $l\gg1$
$\sqrt{l(l+1)}\approx l+\frac12$, относительная неопределённость направления $\sim1/\sqrt l\to0$ —
классический предел.

\begin{figure}[htbp]
\centering
\begin{tikzpicture}[scale=0.9]
  \def\R{2.449} % sqrt(6) for l=2
  \draw[->, thick] (0,-2.9) -- (0,3.2) node[above] {$z$};
  \draw[gray] (0,0) circle (\R);
  \foreach \m in {-2,-1,0,1,2} {
     \draw[dashed, gray] (-3.0,\m) -- (3.0,\m);
     \node[right, font=\small] at (3.05,\m) {$m=\m$};
     \pgfmathsetmacro{\x}{sqrt(6-\m*\m)}
     \draw[very thick, accent, ->] (0,0) -- (\x,\m);
     \draw[accent!50] (-\x,\m) -- (\x,\m);
  }
  \node[font=\small, accent] at (1.9,2.45) {$|\vect L|=\hbar\sqrt6$};
  \draw[keycol, thick] (0,2) ellipse ({sqrt(2)} and 0.22);
  \node[font=\small, align=left] at (-2.9,2.6) {конус\\прецессии};
  \draw[->, thin] (-2.3,2.3) -- (-1.2,2.05);
  \node[font=\small] at (0.35,-2.75) {$0$};
\end{tikzpicture}
\caption{Векторная модель для $l=2$: длина вектора момента $\hbar\sqrt{l(l+1)}=\hbar\sqrt6\approx2{,}45\hbar$,
возможные проекции $L_z=m\hbar$, $m=-2,\dots,2$ (единица по оси $z$ — $\hbar$). Для $m=2$ показан конус,
по которому «распределены» неопределённые компоненты $L_x$, $L_y$.}
\label{fig:9:vector}
\end{figure}

\section{Математические выводы}

\begin{derivation}[label=der:9:comm]{коммутатор $[\op L_x,\op L_y]=\ii\hbar\op L_z$}
Используем $[\op x_i,\op p_j]=\ii\hbar\delta_{ij}$, $[\op x_i,\op x_j]=[\op p_i,\op p_j]=0$ и тождество
$[\op A,\op B\op C]=[\op A,\op B]\op C+\op B[\op A,\op C]$.
\[
  [\op L_x,\op L_y]=[\,y\op p_z-z\op p_y,\ z\op p_x-x\op p_z\,]
  =[y\op p_z,z\op p_x]-[y\op p_z,x\op p_z]-[z\op p_y,z\op p_x]+[z\op p_y,x\op p_z].
\]
Второй и третий коммутаторы равны нулю (в них нет сопряжённых пар). В первом не коммутируют
только $\op p_z$ и $z$: $[y\op p_z,z\op p_x]=y\op p_x[\op p_z,z]=-\ii\hbar\,y\op p_x$. В четвёртом:
$[z\op p_y,x\op p_z]=x\op p_y[z,\op p_z]=\ii\hbar\,x\op p_y$. Итого
\[
  [\op L_x,\op L_y]=\ii\hbar(x\op p_y-y\op p_x)=\ii\hbar\op L_z .
\]
Остальные соотношения~\eqref{eq:9:comm} получаются циклической перестановкой $x\to y\to z\to x$.

\emph{Коммутатор с $\op L^2$.} Например,
\begin{align*}
 [\op L^2,\op L_z]&=[\op L_x^2,\op L_z]+[\op L_y^2,\op L_z]
 =\op L_x[\op L_x,\op L_z]+[\op L_x,\op L_z]\op L_x+\op L_y[\op L_y,\op L_z]+[\op L_y,\op L_z]\op L_y\\
 &=-\ii\hbar(\op L_x\op L_y+\op L_y\op L_x)+\ii\hbar(\op L_y\op L_x+\op L_x\op L_y)=0 .
\end{align*}
\end{derivation}

\begin{derivation}[label=der:9:Lz]{$\op L_z=-\ii\hbar\,\partial/\partial\varphi$}
Рассмотрим функцию $\psi(x,y,z)$ как функцию сферических координат. По правилу дифференцирования
сложной функции
\[
  \pdiff{\psi}{\varphi}=\pdiff{\psi}{x}\pdiff{x}{\varphi}+\pdiff{\psi}{y}\pdiff{y}{\varphi}+\pdiff{\psi}{z}\pdiff{z}{\varphi}.
\]
Из $x=r\sin\theta\cos\varphi$, $y=r\sin\theta\sin\varphi$, $z=r\cos\theta$:
$\pdiff{x}{\varphi}=-r\sin\theta\sin\varphi=-y$, $\pdiff{y}{\varphi}=r\sin\theta\cos\varphi=x$, $\pdiff{z}{\varphi}=0$. Следовательно,
\[
  \pdiff{}{\varphi}=x\pdiff{}{y}-y\pdiff{}{x}\quad\Longrightarrow\quad
  \op L_z=-\ii\hbar\Bigl(x\pdiff{}{y}-y\pdiff{}{x}\Bigr)=-\ii\hbar\pdiff{}{\varphi}.
\]
\end{derivation}

\begin{derivation}[label=der:9:eig]{собственные функции и собственные значения $\op L_z$}
Уравнение $-\ii\hbar\,\dd\Phi/\dd\varphi=L_z\Phi$ — линейное первого порядка:
\[
  \frac{\dd\Phi}{\Phi}=\frac{\ii L_z}{\hbar}\dd\varphi\quad\Longrightarrow\quad\Phi(\varphi)=C\,\ee^{\ii L_z\varphi/\hbar}.
\]
\emph{Условие однозначности:} $\Phi(\varphi+2\pi)=\Phi(\varphi)$ $\Rightarrow$ $\ee^{2\pi\ii L_z/\hbar}=1$ $\Rightarrow$
$L_z/\hbar=m$, $m\in\mathbb Z$. Итак, $L_z=m\hbar$.

\emph{Нормировка:} $\int_0^{2\pi}|C|^2\dd\varphi=2\pi|C|^2=1$ $\Rightarrow$ $C=1/\sqrt{2\pi}$.

\emph{Ортогональность:} $\frac{1}{2\pi}\int_0^{2\pi}\ee^{\ii(m-m')\varphi}\dd\varphi=\delta_{mm'}$.

\emph{Эрмитовость} $\op L_z$ на периодических функциях проверяется интегрированием по частям:
внеинтегральный член $[\Phi_1^*\Phi_2]_0^{2\pi}=0$ из-за периодичности. Поэтому собственные значения
действительны, что согласуется с полученным результатом.
\end{derivation}

\begin{derivation}[label=der:9:Y10]{проверка: $Y_{10}\propto\cos\theta$ — собственная функция $\op L^2$ с $l=1$}
$Y_{10}$ не зависит от $\varphi$, поэтому в~\eqref{eq:9:L2} остаётся только $\theta$-часть:
\[
  \frac{1}{\sin\theta}\diff{}{\theta}\Bigl(\sin\theta\diff{\cos\theta}{\theta}\Bigr)
  =\frac{1}{\sin\theta}\diff{}{\theta}\bigl(-\sin^2\theta\bigr)=\frac{-2\sin\theta\cos\theta}{\sin\theta}=-2\cos\theta .
\]
Значит, $\op L^2Y_{10}=-\hbar^2(-2)Y_{10}=2\hbar^2Y_{10}=\hbar^2\cdot1\cdot(1+1)\,Y_{10}$, т.\,е. $l=1$; и
$\op L_zY_{10}=0$, т.\,е. $m=0$. Аналогично для $Y_{11}\propto\sin\theta\,\ee^{\ii\varphi}$:
$\frac{1}{\sin\theta}\frac{\dd}{\dd\theta}(\sin\theta\cos\theta)-\frac{\sin\theta}{\sin^2\theta}=\frac{\cos2\theta-1}{\sin\theta}=-2\sin\theta$ —
снова $L^2=2\hbar^2$, а $L_z=\hbar$.
\end{derivation}

\begin{derivation}[label=der:9:sum]{сумма $\sum_{m=-l}^{l}m^2$}
$\sum_{m=-l}^{l}m^2=2\sum_{m=1}^{l}m^2=2\cdot\frac{l(l+1)(2l+1)}{6}=\frac{l(l+1)(2l+1)}{3}$ — используется в оценке
$\avg{L_z^2}=\hbar^2l(l+1)/3$ выше.
\end{derivation}

\section{Итоговые формулы}

\keyformula{\op{\vect L}=-\ii\hbar\,\vect r\times\nabla,\qquad [\op L_x,\op L_y]=\ii\hbar\op L_z,\qquad [\op L^2,\op L_i]=0}
 {$\op L_i$ — проекции, $\op L^2$ — квадрат момента; циклические перестановки $x\to y\to z$.}
 {орбитальный момент частицы (те же коммутаторы — для любого момента, включая спин).}
 {какие величины одновременно измеримы ($L^2$ и одна проекция).}

\keyformula{\op L_z=-\ii\hbar\pdiff{}{\varphi},\qquad \Phi_m=\frac{\ee^{\ii m\varphi}}{\sqrt{2\pi}},\qquad L_z=m\hbar,\ m\in\mathbb Z}
 {$\varphi$ — азимутальный угол, $m$ — магнитное квантовое число.}
 {однозначная волновая функция (орбитальный момент).}
 {спектр и собственные функции проекции момента.}

\keyformula{L^2=\hbar^2l(l+1),\qquad l=0,1,2,\dots;\qquad m=-l,\dots,l\ \ (2l+1\ \text{значений})}
 {$l$ — орбитальное квантовое число; $Y_{lm}(\theta,\varphi)$ — сферические функции.}
 {без вывода; орбитальный момент.}
 {классификация состояний в центральном поле ($s,p,d,f$); кратность вырождения.}

\section{Физическая интерпретация}

\textbf{Почему квантуется момент.} Требование однозначности $\Psi$ при повороте на $2\pi$ приводит
к $L_z=m\hbar$; дискретность $L^2$ — следствие условий конечности на сфере. Бор постулировал
$L=n\hbar$; квантовая механика даёт другое: $|\vect L|=\hbar\sqrt{l(l+1)}$, причём $l=0$ возможно
(основное состояние водорода имеет нулевой момент).

\textbf{Сохранение момента.} В центральном поле $U(r)$ операторы $\op L^2$ и $\op L_z$ коммутируют с
$\op H$, поэтому их собственные значения — «хорошие» квантовые числа. Энергия не зависит от $m$
(нет выделенного направления) — вырождение кратности $2l+1$. Магнитное поле выделяет ось $z$ и
снимает это вырождение (эффект Зеемана, \tref{t:12}) — отсюда название «магнитное» число.

\textbf{Связь с экспериментом.} Пространственное квантование проекции момента наблюдается в опыте
Штерна--Герлаха и в расщеплении спектральных линий в магнитном поле на $2l+1$ (или $2j+1$) компонент.

\section{Что необходимо знать на экзамене}

\begin{examcheck}
\checkhead{Определения}
\begin{itemize}
 \item $\op{\vect L}=\vect r\times\op{\vect p}$, компоненты в декартовых координатах; квантовые числа $l$, $m$; обозначения $s,p,d,f$.
\end{itemize}
\checkhead{Формулы}
\begin{itemize}
 \item $[\op L_x,\op L_y]=\ii\hbar\op L_z$ (и циклические), $[\op L^2,\op L_i]=0$;
 \item $\op L_z=-\ii\hbar\,\partial_\varphi$, $\Phi_m=\ee^{\ii m\varphi}/\sqrt{2\pi}$, $L_z=m\hbar$;
 \item $L^2=\hbar^2l(l+1)$, $|m|\le l$, кратность $2l+1$ (без вывода); $Y_{00}$, $Y_{10}$.
\end{itemize}
\checkhead{Выводы}
\begin{itemize}
 \item коммутатор $[\op L_x,\op L_y]$; переход к $\op L_z=-\ii\hbar\,\partial_\varphi$;
 \item решение задачи на собственные значения $\op L_z$ с условием однозначности и нормировкой.
\end{itemize}
\checkhead{Интерпретации}
\begin{itemize}
 \item почему одновременно измеримы лишь $L^2$ и одна проекция; почему $|m|\le l$ и $L>l\hbar$;
 \item векторная модель; классический предел $l\gg1$; связь с сохранением момента и вырождением по $m$.
\end{itemize}
\end{examcheck}
